% ── Parties Taxonomy ─────────────────────────────────────────────────────────
% Edit this file to customise the "Parties Taxonomy" section.

The party catalogue records metadata for every party, electoral list, or
individual candidate that appears in at least one election result file in
BLUE\_DB.

\paragraph{Diversity of entries.}
The \texttt{parties/parties.csv} file contains several distinct types of
entry.  Most rows represent organised parties or standing lists.  Rows where
\bluevar{coalition\_members} is non-empty represent electoral coalitions whose
member parties also appear as separate rows.  Rows where the field
\bluevar{individual\_candidate} is \texttt{Y} represent independent
candidates or candidate groups without a formal party structure.

A special case arises for French elections, where the Ministry of the Interior
assigns many small-party or independent candidates a \emph{nuance politique} (political shading) drawn from a
fixed set of administrative labels rather than the name of a registered party.
These nuances are included as distinct entries in the parties file (e.g.\
\texttt{DVD} for \emph{divers droite}). However, they do not correspond to
specific legal party entities.

\subsection{European Parliament groups}

Parties are classified into the European Parliament (EP) political group of
their MEPs, or of the associated European-level party.  Because parties may
switch groups across parliamentary terms, the \bluevar{ep\_group} field stores
a JSON array of membership spells.  BLUE\_DB aggregates groups that share a
direct institutional lineage (e.g.\ the \emph{Party of European Socialists}
and its successor \emph{S\&D}) into a single row in
\texttt{parties/ep\_groups.csv}.  The \textbf{Parties} column below
counts the number of national parties affiliated with each group.

\VAR{ep_groups_table}

\subsection{European parties}

The \bluevar{eu\_party} field records the European-level party (federation) to
which a national party is affiliated.  BLUE\_DB considers only European
political parties formally registered with the Authority for European
Political Parties and European Political Foundations (APPF).  Some European
parties have undergone name changes or mergers; former names are recorded in
the \bluevar{other\_names} field of \texttt{parties/eu\_parties.csv}.
The \textbf{Parties} column below counts the number of national parties that
have held membership.

\VAR{european_parties_table}

\subsection{Ideology}

The \bluevar{ideology} field assigns each party to a broad ideological family.
The classification follows the EP group or European party affiliation where
one is available, using the defaults in the table below.

\begin{longtable}{p{6.5cm}p{4cm}p{4cm}}
\toprule
\textbf{Ideology} & \textbf{EP groups} & \textbf{European parties} \\
\midrule
\endhead
\bottomrule
\endlastfoot
center, center-right or moderate right
  & \textbf{EPP Group}, \textbf{Renew}
  & \textbf{EPP}, \textbf{ALDE}, \textbf{EDP} \\
left or center-left, incl.\ non-radical greens
  & \textbf{S\&D}, \textbf{Greens/EFA}
  & \textbf{PES}, \textbf{EGP} \\
radical left
  & \textbf{GUE/NGL}
  & \textbf{PEL}, \textbf{ELA} \\
radical right
  & \textbf{PfE}, \textbf{ESN Group}, \textbf{UEN}
  & \textbf{Patriots.eu}, \textbf{ESN}, \textbf{ECR}, \textbf{APF}, \textbf{AENM} \\
individually assessed
  & \textbf{ECR Group}, \textbf{EFDD}, \textbf{NI}
  & \textbf{EFA}, \textbf{ECPP} \\
\end{longtable}
For parties whose group membership falls under \emph{individually assessed} (ECR Group,
EFDD, NI), as well as for smaller parties, regional lists, and independent
candidates not connected to a European-level structure, the classification
relies on manual evaluation drawing primarily on Wikipedia sources.


\VAR{ideology_table}
